% Traduction française COMIMa de l'énoncé original (anglais).
On se donne une droite $a$ et un point $P$. En traçant le moins de lignes possible (cercles ou droites), construire la perpendiculaire à $a$ passant par $P$. Si $P \notin a$, la construction est bien connue des lycéens. Mais supposons $P \in a$ : la construction minimale est peu connue. Voir la solution dans l'ouvrage pour la preuve de cette affirmation.
